\cleardoublepage
\begin{thebibliography}{99}
\addcontentsline{toc}{chapter}{Bibliografía}

\bibitem{zeppenfeld2022}
Zeppenfeld, K. y colaboradores.
\newblock 2022 ESC Guidelines for the management of patients with ventricular arrhythmias and the prevention of sudden cardiac death.
\newblock \emph{European Heart Journal}, 43(40), 3997--4126, 2022.
\newblock \url{https://doi.org/10.1093/eurheartj/ehac262}.

\bibitem{kligfield2007}
Kligfield, P., Gettes, L. S., Bailey, J. J. y colaboradores.
\newblock Recommendations for the standardization and interpretation of the electrocardiogram. Part I. The electrocardiogram and its technology.
\newblock \emph{Circulation}, 115(10), 1306--1324, 2007.
\newblock \url{https://doi.org/10.1161/CIRCULATIONAHA.106.180200}.

\bibitem{surawicz2009}
Surawicz, B., Childers, R., Deal, B. J., Gettes, L. S. y colaboradores.
\newblock AHA/ACCF/HRS recommendations for the standardization and interpretation of the electrocardiogram. Part III: Intraventricular conduction disturbances.
\newblock \emph{Circulation}, 119(10), e235--e240, 2009.
\newblock \url{https://doi.org/10.1161/CIRCULATIONAHA.108.191095}.

\bibitem{rautaharju2009st}
Rautaharju, P. M., Surawicz, B. y Gettes, L. S.
\newblock AHA/ACCF/HRS recommendations for the standardization and interpretation of the electrocardiogram. Part IV: The ST segment, T and U waves, and the QT interval.
\newblock \emph{Circulation}, 119(10), e241--e250, 2009.
\newblock \url{https://doi.org/10.1161/CIRCULATIONAHA.108.191096}.

\bibitem{wilde2023scbt}
Wilde, A. A. M., Amin, A. S. y colaboradores.
\newblock Use, misuse, and pitfalls of the drug challenge test in the diagnosis of the Brugada syndrome.
\newblock \emph{European Heart Journal}, 44(27), 2427--2439, 2023.
\newblock \url{https://doi.org/10.1093/eurheartj/ehad295}.

\bibitem{ecgpediaType1Image}
CardioNetworks y ECGpedia.
\newblock Brugada syndrome type1 example4.
\newblock Wikimedia Commons, licencia CC BY-SA 3.0.
\newblock \url{https://commons.wikimedia.org/wiki/File:Brugada_syndrome_type1_example4_(CardioNetworks_ECGpedia).png}.

\bibitem{napolitano2006}
Napolitano, C. y Priori, S. G.
\newblock Brugada syndrome.
\newblock \emph{Orphanet Journal of Rare Diseases}, 1, 35, 2006.
\newblock \url{https://doi.org/10.1186/1750-1172-1-35}.

\bibitem{ribeiro2020}
Ribeiro, A. H. y colaboradores.
\newblock Automatic diagnosis of the 12-lead ECG using a deep neural network.
\newblock \emph{Nature Communications}, 11, 1760, 2020.
\newblock \url{https://doi.org/10.1038/s41467-020-15432-4}.

\bibitem{wagner2020}
Wagner, P., Strodthoff, N., Bousseljot, R. D. y colaboradores.
\newblock PTB-XL, a large publicly available electrocardiography dataset.
\newblock \emph{Scientific Data}, 7, 154, 2020.
\newblock \url{https://doi.org/10.1038/s41597-020-0495-6}.

\bibitem{saleh2026brugada}
Saleh, K. y colaboradores.
\newblock AI-ECG classification for Brugada syndrome: A study of machine learning techniques to optimise for limited datasets.
\newblock \emph{PLOS Digital Health}, 5(2), e0001222, 2026.
\newblock \url{https://doi.org/10.1371/journal.pdig.0001222}.

\bibitem{gliner2020}
Gliner, V., Keidar, N., Makarov, V. y colaboradores.
\newblock Automatic classification of healthy and disease conditions from images or digital standard 12-lead electrocardiograms.
\newblock \emph{Scientific Reports}, 10, 16331, 2020.
\newblock \url{https://doi.org/10.1038/s41598-020-73060-w}.

\bibitem{sangha2022}
Sangha, V. y colaboradores.
\newblock Automated multilabel diagnosis on electrocardiographic images and signals.
\newblock \emph{Nature Communications}, 13, 1583, 2022.
\newblock \url{https://doi.org/10.1038/s41467-022-29153-3}.

\bibitem{gliner2023}
Gliner, V., Makarov, V., Avetisyan, A. I., Schuster, A. y Yaniv, Y.
\newblock Using domain adaptation for classification of healthy and disease conditions from mobile-captured images of standard 12-lead electrocardiograms.
\newblock \emph{Scientific Reports}, 13, 14023, 2023.
\newblock \url{https://doi.org/10.1038/s41598-023-40693-6}.

\bibitem{sayadi2010}
Sayadi, O., Shamsollahi, M. B. y Clifford, G. D.
\newblock Synthetic ECG generation and Bayesian filtering using a Gaussian wave-based dynamical model.
\newblock \emph{Physiological Measurement}, 31(10), 1309--1329, 2010.
\newblock \url{https://doi.org/10.1088/0967-3334/31/10/002}.

\bibitem{liu2026}
Liu, R., Bai, Y., Yue, X. y Zhang, P.
\newblock Teaching multimodal LLMs to comprehend 12-lead electrocardiographic images.
\newblock \emph{npj Digital Medicine}, 9, 349, 2026.
\newblock \url{https://doi.org/10.1038/s41746-026-02551-3}.

\bibitem{dong2024survey}
Dong, Q., Li, L., Dai, D., Zheng, C. y colaboradores.
\newblock A survey on in-context learning.
\newblock \emph{Proceedings of the 2024 Conference on Empirical Methods in Natural Language Processing}, 1107--1128, 2024.
\newblock \url{https://aclanthology.org/2024.emnlp-main.64/}.

\bibitem{alayrac2022}
Alayrac, J. B., Donahue, J., Luc, P. y colaboradores.
\newblock Flamingo: A visual language model for few-shot learning.
\newblock \emph{Advances in Neural Information Processing Systems}, 35, 23716--23736, 2022.
\newblock \url{https://doi.org/10.52202/068431-1723}.

\bibitem{camba2026icl}
Camba, S., Otero, A., García, D., Sánchez, L. y Costa, N.
\newblock Can In-Context Learning Enable Large Vision Language Models to Detect ECG Abnormalities?
\newblock En \emph{Artificial Intelligence in Biomedicine. CIABiomed 2025}, Lecture Notes in Computer Science, vol. 16148, 200--213. Springer, Cham, 2026.
\newblock \url{https://doi.org/10.1007/978-3-032-10661-2_16}.

\bibitem{costa2026fewshot}
Costa, N., Otero, A., García Iglesias, D. y Sánchez, L.
\newblock Few-shot ECG analysis with large vision language models in data-scarce clinical settings.
\newblock \emph{Computers in Biology and Medicine}, 213, 111816, 2026.
\newblock \url{https://doi.org/10.1016/j.compbiomed.2026.111816}.

\bibitem{gemma4card}
Google DeepMind.
\newblock Gemma 4 model card.
\newblock Google AI for Developers.
\newblock \url{https://ai.google.dev/gemma/docs/core/model_card_4}.

\bibitem{he2016}
He, K., Zhang, X., Ren, S. y Sun, J.
\newblock Deep residual learning for image recognition.
\newblock \emph{Proceedings of the IEEE Conference on Computer Vision and Pattern Recognition}, 770--778, 2016.
\newblock \url{https://doi.org/10.1109/CVPR.2016.90}.

\bibitem{selvaraju2017}
Selvaraju, R. R. y colaboradores.
\newblock Grad-CAM: Visual explanations from deep networks via gradient-based localization.
\newblock \emph{Proceedings of the IEEE International Conference on Computer Vision}, 618--626, 2017.
\newblock \url{https://doi.org/10.1109/ICCV.2017.74}.

\bibitem{ganin2016}
Ganin, Y. y colaboradores.
\newblock Domain-adversarial training of neural networks.
\newblock \emph{Journal of Machine Learning Research}, 17(59), 1--35, 2016.
\newblock \url{https://jmlr.org/papers/v17/15-239.html}.

\bibitem{chen2020simclr}
Chen, T., Kornblith, S., Norouzi, M. y Hinton, G.
\newblock A simple framework for contrastive learning of visual representations.
\newblock \emph{Proceedings of the 37th International Conference on Machine Learning}, PMLR 119, 1597--1607, 2020.
\newblock \url{https://proceedings.mlr.press/v119/chen20j.html}.

\bibitem{sun2016deepcoral}
Sun, B. y Saenko, K.
\newblock Deep CORAL: Correlation alignment for deep domain adaptation.
\newblock En \emph{European Conference on Computer Vision Workshops}, 443--450, 2016.
\newblock \url{https://doi.org/10.1007/978-3-319-49409-8_35}.

\bibitem{long2015dan}
Long, M., Cao, Y., Wang, J. y Jordan, M. I.
\newblock Learning transferable features with deep adaptation networks.
\newblock \emph{Proceedings of the 32nd International Conference on Machine Learning}, PMLR 37, 97--105, 2015.
\newblock \url{https://proceedings.mlr.press/v37/long15.html}.

\bibitem{loshchilov2019}
Loshchilov, I. y Hutter, F.
\newblock Decoupled weight decay regularization.
\newblock \emph{International Conference on Learning Representations}, 2019.
\newblock \url{https://arxiv.org/abs/1711.05101}.

\par\smallskip
\bibitem{costacortez2026brugada}
Costa Cortez, N. y García Iglesias, D.
\newblock Brugada-HUCA: 12-Lead ECG Recordings for the Study of Brugada Syndrome. Versión 1.0.0.
\newblock PhysioNet, 2026.
\newblock \url{https://doi.org/10.13026/0m2w-dy83}.

\bibitem{strodthoff2021ptbxl}
Strodthoff, N., Wagner, P., Schaeffter, T. y Samek, W.
\newblock Deep learning for ECG analysis: Benchmarks and insights from PTB-XL.
\newblock \emph{IEEE Journal of Biomedical and Health Informatics}, 25(5), 1519--1528, 2021.
\newblock \url{https://doi.org/10.1109/JBHI.2020.3022989}.

\bibitem{gow2023mimicecg}
Gow, B., Pollard, T., Nathanson, L. A. y colaboradores.
\newblock MIMIC-IV-ECG: Diagnostic electrocardiogram matched subset. Versión 1.0.
\newblock PhysioNet, 2023.
\newblock \url{https://doi.org/10.13026/4nqg-sb35}.

\bibitem{mckeen2025ecgfm}
McKeen, K. y colaboradores.
\newblock ECG-FM: An open electrocardiogram foundation model.
\newblock \emph{JAMIA Open}, 8(5), ooaf122, 2025.
\newblock \url{https://doi.org/10.1093/jamiaopen/ooaf122}.

\bibitem{incartdb}
St.~Petersburg Institute of Cardiological Technics.
\newblock St Petersburg INCART 12-lead arrhythmia database. Versión 1.0.0.
\newblock PhysioNet, 2008.
\newblock \url{https://physionet.org/content/incartdb/1.0.0/}.

\bibitem{he2018pyramid}
He, J., Sun, L., Rong, J., Wang, H. y Zhang, Y.
\newblock A pyramid-like model for heartbeat classification from ECG recordings.
\newblock \emph{PLOS ONE}, 13(11), e0206593, 2018.
\newblock \url{https://doi.org/10.1371/journal.pone.0206593}.

\bibitem{reyna2024challenge}
Reyna, M. A., Deepanshi, Weigle, J. T. y colaboradores.
\newblock Digitization and classification of ECG images: The George B. Moody PhysioNet Challenge 2024.
\newblock \emph{Computing in Cardiology}, 51, 1--4, 2024.
\newblock \url{https://doi.org/10.22489/CinC.2024.011}.

\bibitem{uniovi2026guiatfe}
Facultad de Formación del Profesorado y Educación, Universidad de Oviedo.
\newblock Guía de uso de la inteligencia artificial en los TFE. DOC-SGIC-UO-FFPE-35, 2026.
\newblock \href{https://digibuo.uniovi.es/dspace/bitstream/handle/10651/83693/Guia_IA_TFE.pdf?sequence=1\&isAllowed=y}{Disponible en DigiBUO}.

\bibitem{openaicodex}
OpenAI.
\newblock Codex [asistente de programación].
\newblock \url{https://openai.com/codex/}.

\bibitem{perezalday2020challenge}
Perez Alday, E. A., Gu, A., Shah, A. J. y colaboradores.
\newblock Classification of 12-lead ECGs: the PhysioNet/Computing in Cardiology Challenge 2020.
\newblock \emph{Physiological Measurement}, 41(12), 124003, 2020.
\newblock \url{https://doi.org/10.1088/1361-6579/abc960}.

\bibitem{qtdb}
Laguna, P., Mark, R. G., Goldberger, A. L. y Moody, G. B.
\newblock A database for evaluation of algorithms for measurement of QT and other waveform intervals in the ECG.
\newblock \emph{Computers in Cardiology}, 24, 673--676, 1997.
\newblock \url{https://physionet.org/content/qtdb/1.0.0/}.

\bibitem{jimenezperez2021delineation}
Jimenez-Perez, G., Alcaine, A. y Camara, O.
\newblock Delineation of the electrocardiogram with a mixed-quality-annotations dataset using convolutional neural networks.
\newblock \emph{Scientific Reports}, 11, 863, 2021.
\newblock \url{https://doi.org/10.1038/s41598-020-79512-7}.

\bibitem{calburean2024brugada}
Călburean, P.-A., Pannone, L., Monaco, C. y colaboradores.
\newblock Predicting and recognizing drug-induced type I Brugada pattern using ECG-based deep learning.
\newblock \emph{Journal of the American Heart Association}, 13(10), e033148, 2024.
\newblock \url{https://doi.org/10.1161/JAHA.123.033148}.

\bibitem{mitbih}
Moody, G. B. y Mark, R. G.
\newblock MIT-BIH Arrhythmia Database. Versión 1.0.0.
\newblock PhysioNet, 2005.
\newblock \url{https://doi.org/10.13026/C2F305}.

\bibitem{reyna2022challenge}
Reyna, M., Sadr, N., Gu, A. y colaboradores.
\newblock Will Two Do? Varying Dimensions in Electrocardiography: The PhysioNet/Computing in Cardiology Challenge 2021. Versión 1.0.3.
\newblock PhysioNet, 2022.
\newblock \url{https://doi.org/10.13026/34va-7q14}.

\bibitem{russakovsky2015imagenet}
Russakovsky, O., Deng, J., Su, H. y colaboradores.
\newblock ImageNet Large Scale Visual Recognition Challenge.
\newblock \emph{International Journal of Computer Vision}, 115(3), 211--252, 2015.
\newblock \url{https://doi.org/10.1007/s11263-015-0816-y}.

\bibitem{torchvisionresnet18}
PyTorch.
\newblock \texttt{resnet18}: Arquitectura y pesos preentrenados en ImageNet-1K.
\newblock Documentación de Torchvision.
\newblock \href{https://docs.pytorch.org/vision/stable/models/generated/torchvision.models.resnet18.html}{Documentación de ResNet18}.

\bibitem{kwon2023vllm}
Kwon, W., Li, Z., Zhuang, S. y colaboradores.
\newblock Efficient Memory Management for Large Language Model Serving with PagedAttention.
\newblock \emph{Proceedings of the 29th ACM Symposium on Operating Systems Principles}, 611--626, 2023.
\newblock \url{https://doi.org/10.1145/3600006.3613165}.

\bibitem{sklearnaverageprecision}
Scikit-learn.
\newblock \texttt{average\_precision\_score}.
\newblock Documentación de las métricas de clasificación.
\newblock \href{https://scikit-learn.org/stable/modules/generated/sklearn.metrics.average_precision_score.html}{Documentación de la precisión media}.

\end{thebibliography}
